\section{Applications to Training and Education}
\label{sec:9}

\subsection{Medical communication and debriefing}

A medical extension will require educator-defined communication events and a scenario-specific intervention policy, supported by TEAM validation, nontechnical-skills systems, and physiological analysis of simulation ~\cite{p66,p4,p67,p69}. The proposed comparison distinguishes immediate assistance from delayed reflection where educators judge both appropriate. Debriefing literature supports examining reflective feedback, while the learning-performance distinction requires a new unassisted scenario before claiming skill retention ~\cite{p9,p62}. The agent's target is trainee teamwork; it will not independently select patient treatment.

\missing{[TO COMPLETE: medical partner, participant expertise, scenario and roles, critical-event rubric, validated outcomes, instructor override, delayed or debriefing condition, unassisted transfer scenario, and results.]}

\subsection{Peer explanation and independent understanding}

An education extension will target explanation and shared monitoring rather than answer delivery, drawing on socially shared regulation and pedagogical scaffolding ~\cite{p29,p27,p51}. An individual assessment before and after collaboration, followed by an unassisted new problem, is proposed to test whether team support benefits each learner. Productive-failure and learning-performance research motivate preserving useful struggle and separating supported completion from learning ~\cite{p31,p62}.

\missing{[TO COMPLETE: learning objectives, concept inventory or scoring rubric, prior-knowledge assessment, group task, scaffold withdrawal rule, delayed interval, unassisted transfer problem, and results.]}

\subsection{What can generalize}

The shared object across settings is an evidence-to-decision process, not a universal physiological threshold. Medical communication measures and educational regulation frameworks define different objectives ~\cite{p66,p29}. We propose evaluating transfer of feature representations, event definitions, and policy constraints separately; success in the primary collaboration task will not establish effectiveness in either extension.

\missing{[TO COMPLETE: which extension was actually evaluated; transferred versus re-trained components; domain-specific changes; and evidence-supported limits of generalization.]}
